\section{An alternative source: an explicit equational proof kernel}
\label{sec:relational-kernel}

This appendix supplies a concrete local component of a different source
encoding.  It does not supply a new consistency equation.  The broad route
through finitely axiomatized class theory and relation algebra is prior
work: Carl and Moroz used finite NBG in their Diophantine provability
construction, and Formisano, Gavazzi, and Omodeo subsequently proposed a
Tarski--Givant relational version.  The latter paper does not print its
proper-axiom tuple; its Section~3.3 leaves those formulas outside its scope.
Consequently it cannot, by citation alone, calibrate a concrete new
polynomial to ZFC.  Our contribution here is an independently derived,
completely explicit local inference polynomial and an executable finite
proof checker for a supplied finite collection of equation schemata.
No priority claim is made for the local algebraic identity.

\subsection{A total coding of terms}

Work over $\mathbb N=\{0,1,2,\ldots\}$.  Let the abstract term language have
two constants $\epsilon,\iota$, two unary constructors $N,V$, and two binary
constructors $U,M$.  These will eventually stand for membership, identity,
complement, converse, union, and relation composition, respectively.  The
coding theorem itself is purely syntactic and does not require that
interpretation.  Set
\begin{align*}
 [\epsilon]&=0,& [\iota]&=1,\\
 [N(A)]&=4[A]+2,&[V(A)]&=4[A]+3,\\
 B(a,z)&=2(a+z)(a+z+1)+4z+4,\\
 [U(A,Z)]&=B([A],[Z]),&[M(A,Z)]&=B([A],[Z])+1.
\end{align*}
Every displayed constructor is a polynomial with integer coefficients.
In particular there is no division in the code of a compound term.
This is the polynomial form of the numbering in Section~3.1 of
Formisano--Gavazzi--Omodeo; the complete bijection proof is included here
to make the component independently checkable.

\paragraph{Theorem A.1 (total coding).}
The coding map is a bijection from these finite terms onto $\mathbb N$.

\paragraph{Proof.}
The Cantor function
\[
 C(a,z)=\frac{(a+z)(a+z+1)}2+z
\]
is a bijection $\mathbb N^2\longrightarrow\mathbb N$.  Indeed, for a given
$g$, the unique $s$ satisfying
$s(s+1)/2\leq g<(s+1)(s+2)/2$ gives
$z=g-s(s+1)/2$ and $a=s-z$.  Then $0\leq z\leq s$, so these are natural
numbers, and the formulas also prove uniqueness.  Moreover
$C(a,z)\geq\max(a,z)$.

The two constants occupy $0,1$.  Codes congruent to $2$ and $3$ modulo $4$
have the unique forms $4a+2$ and $4a+3$, respectively.  The remaining
numbers are $4(g+1)$ and $4(g+1)+1$ with $g\geq0$; the unique inverse of
$C$ identifies their two children.  Since $B=4(C+1)$, these are precisely
the two binary constructor cases.  In every case the child codes are
strictly smaller than the parent code.  Strong induction on a natural
number therefore constructs one finite term having that code; the
disjoint residue cases and the uniqueness of the Cantor inverse prove
that this term is unique.  This also proves injectivity of the original
coding map.  All inductions and integer computations are available in
ZFC. \hfill$\square$

Thus a natural-number witness never needs a separate certificate that
it denotes a valid term.  Equalities are stored as two coordinates
$(a,b)$, representing the equation between the uniquely decoded terms.
This avoids packing and immediately unpacking an equality as a further
Cantor pair.

\subsection{The seven-variable local polynomial}

Let $(a,b)$ and $(c,d)$ be two already proved equations and let $(u,v)$
be a proposed conclusion.  Introduce only one additional natural witness
$z$.  Define
\begin{align*}
 T&=(a-c)^2+(u-b)^2+(v-d)^2,\\
 q&=u-4a-2,\\
 A&=\bigl(q(q-1)\bigr)^2+\bigl(v-u-4(b-a)\bigr)^2,\\
 r&=u-B(a,z),\\
 D&=\bigl(r(r-1)\bigr)^2
       +\bigl(v-u-2(b-a)(a+b+2z+1)\bigr)^2,\\
 K(u,v,a,b,c,d,z)&=TAD.
\end{align*}
There are seven variables in $K$, not seven witnesses for the whole
provability problem: six coordinates are inputs and outputs of a single
local inference.  The history encoding is a separate cost.

\paragraph{Theorem A.2 (exact local inference).}
For natural values of its seven arguments, $K=0$ holds if and only if at
least one of the following five inference cases holds:
\begin{enumerate}
\item $a=c$, $u=b$, and $v=d$ (common-left transitivity);
\item $u=4a+2$ and $v=4b+2$ (complement congruence);
\item $u=4a+3$ and $v=4b+3$ (converse congruence);
\item $u=B(a,z)$ and $v=B(b,z)$ (union in the right argument);
\item $u=B(a,z)+1$ and $v=B(b,z)+1$ (composition in the right argument).
\end{enumerate}
Here ``in the right argument'' means adjoining the same right operand
$z$ to both members of a proved equality; it does not mean changing that
right operand.

\paragraph{Proof.}
The integer ring has no zero divisors, so $K=0$ is equivalent to the
vanishing of one factor.  Each factor is a sum of squares of integers.
For $T$, this gives exactly case~1.  For $A$, it gives
$q(q-1)=0$ and $v-u=4(b-a)$.  Hence $q$ is $0$ or $1$ and the same unary
tag is used on both sides, proving cases~2 and~3.

For $D$, it gives $r\in\{0,1\}$ and the stated difference equation.
Direct expansion gives the polynomial identity
\[
 B(b,z)-B(a,z)=2(b-a)(a+b+2z+1).
\]
Thus $u=B(a,z)+r$ and $v=B(b,z)+r$, proving cases~4 and~5.
Every listed case conversely annihilates its indicated factor.
The argument establishing this union of arithmetic cases works over all
integers; nonnegativity is used only for the term interpretation from
Theorem~A.1. \hfill$\square$

The two adjacent constructor tags are therefore selected without
introducing a Boolean variable.  The finite-difference identity also
avoids calculating $B(a,z)$ and $B(b,z)$ independently.  The included
straight-line circuit has 41 live binary additions, subtractions, and
multiplications; a square counts as one multiplication, and its constants
are $1,2,4$.  The circuit is an evaluator for this local polynomial, not
an operation count for a completed consistency equation.

\paragraph{Exact degree.}
The three factors have degrees $2,4,8$, respectively, so
$\deg K\leq14$.  Specializing $a=t$ and all other inputs to zero gives
leading terms $t^2$, $256t^4$, and $16t^8$.  Therefore the resulting
univariate polynomial has leading term $4096t^{14}$, proving
$\deg K=14$.

\subsection{Finite proof checking and the remaining obligations}

Fix a finite list of equation schemata.  Each side is a finite term
template over the six constructors, with a finite number of schematic
variables.  Supplying natural values for those variables produces an
exact instance by the constructor polynomials.  The code contains no
axiomhood oracle: the list of templates is literal input data.

A proof line records its conclusion, and either (i) a schema index and
its instantiation values, or (ii) two strictly earlier line indices and
the witness $z$.  A line of kind~(ii) is accepted exactly when $K=0$ for
its two parent equations.  One-parent congruence steps repeat the same
parent index.  The first line must be a schema instance.  This explicit
requirement excludes an unjustified cyclic first line.

\paragraph{Theorem A.3 (finite-checker adequacy).}
For the supplied schema list, the checker accepts a proof of a target
equation exactly when that equation has a finite derivation from those
schemata using the five inference cases in Theorem~A.2.

\paragraph{Proof.}
If the checker accepts, induction over its finite list proves each line.
At a schema line this follows by exact instantiation.  At an inference
line both parents have smaller indices, so the induction hypothesis and
Theorem~A.2 give its conclusion.  Conversely, list the nodes of a finite
derivation tree in a children-before-parent order.  Record the two
earlier indices for a transitivity step, duplicate the earlier index
for a congruence step, and use the code of its right operand for a binary
congruence.  For the other steps $z=0$ suffices.  Theorem~A.2 makes the
kernel vanish at each inference, and exact constructor evaluation
verifies every schema instance. \hfill$\square$

This theorem concerns the stated finite equational calculus, without a
claim that an arbitrary schema list expresses set theory.  A finished
replacement for the existing Con(ZFC) equation still requires all of:
\begin{enumerate}
\item A completely printed finite relational class-theory axiom tuple,
including the logical relation-algebra schemata and a precise target
equation expressing collapse.
\item Proof transformations, formalizable in ZFC, establishing that this
target is derivable exactly when the chosen ZFC calculus is inconsistent.
In particular the quasi-projection assumptions required for the
Tarski--Givant translation must be proved or included correctly.
\item An explicit polynomial compiler for arbitrarily long accepted
finite histories, with every bound, chronology condition, coefficient,
and witness accounted for.
\item Complete global variable, degree, coefficient, and evaluator counts
after that compiler is instantiated.
\end{enumerate}
None of these global costs can be inferred from the seven local
coordinates or 41 local operations.  The component provides a distinct
testable source direction, while leaving the established consistency
equation unchanged.

\subsection{An explicit combined row polynomial}

There is also a direct polynomial form of schema admission.  Suppose
the $j$th supplied schema has at most $m$ schematic variables.  Evaluating
its left and right templates using the constructor polynomials gives
fixed polynomials $L_j(\mathbf x),R_j(\mathbf x)\in
\mathbb Z[x_1,\ldots,x_m]$; unused coordinates are simply ignored.
Define
\[
 H(u,v,\mathbf x)=\prod_{j=1}^{s}
 \bigl((u-L_j(\mathbf x))^2+(v-R_j(\mathbf x))^2\bigr),
 \qquad
 F=K(u,v,a,b,c,d,z)H(u,v,\mathbf x).
\]
The empty product is $1$.  Over natural inputs, $F=0$ means exactly
that the output is a listed schema instance or follows from the two
parents by one of the five inference cases.  The same $m$ schema
coordinates may be shared between all $s$ factors: a zero product
requires only one schema factor to vanish.  Conversely witnesses for
that chosen schema can be padded with zeros.  This observation does
not share witnesses between multiple simultaneously required premises.

If the chosen logical basis includes reflexivity, histories can start
with the fixed justified row $(0,0)$.  Every later row can then carry two
strictly earlier parent indices, including axiom rows, which use the
initial row as harmless dummy parents.  Using $F=0$ for those later rows
avoids a separate Boolean axiom/inference tag.  The first row is fixed
and verified; it is never admitted merely by a self-referential
inference.  Induction gives the same soundness and completeness proof
as above.  This combined row polynomial still has an unbounded-history
compilation obligation.  In particular the variable count $7+m$ is a
local row count, not a variable bound for a Con(ZFC) equation.

\paragraph{Reproduction.}
Run \texttt{python3 relational\_kernel.py}.  The included finite checks
cover 10,000 complete code roundtrips, 10,000 Cantor inverse roundtrips,
2,048 binary constructor cases, 20,000 seeded arithmetic comparisons,
8,640 positive inference cases, and valid/invalid finite proofs.  These
checks supplement the universal proofs above; they do not establish the
missing NBG calibration or global history compilation.

\paragraph{Primary references.}
M.~Carl and B.~Z.~Moroz, \emph{On a Diophantine representation of the
predicate of provability}, MPI preprint~2006/129,
\url{https://webdoc.sub.gwdg.de/ebook/serien/e/mpi_mathematik/2006/129.pdf}.
A.~Formisano, I.~Gavazzi, and E.~G.~Omodeo,
\emph{Diophantine Modeling of Provability in Algebraic Logic}, CILC~2025,
\url{https://ceur-ws.org/Vol-4003/paper13.pdf}.
